In this section, we describe the empirical evaluation datasets, baselines, ablations, and metrics, and then discuss the main numerical results.

\subsection{Datasets}
\label{sec:data}
Evaluation data were obtained from the open-access PK-DB repository, which aggregates pharmacokinetic measurements from clinical and pre-clinical studies~\citep{grzegorzewski2021pkdb}.
%
From PK-DB, we curated plasma concentration--time measurements from a Phase~I study in healthy volunteers investigating the safety and dosing of a combination of drugs and their metabolites~\citep{lenuzza2016safety}.
%
This dataset comprises 18 compounds, including parent drugs and primary metabolites, with 2019 valid concentration--time observations.
%
Sampling frequencies vary substantially both within and across compounds, resulting in sparse and irregularly sampled time series; each compound is measured in 10 individuals, with up to 15 observations per individual.
%
We also include two pharmacokinetic datasets on indometacin and theophylline (\texttt{Indometh} and \texttt{Theoph} from the R \texttt{datasets} package)~\citep{Rsoftware}.
%
Together, these 20 empirical compound-level datasets reflect the small-cohort and sparsely sampled regime targeted by PFF.

\subsection{Baselines and Ablations}
\label{sec:baselines}

We compare Prior-Fitted Functional Flows (PFF) against one classical pharmacometric baseline and three neural baselines: nonlinear mixed-effects modeling (NLME), NODE-PK, Transformer, and \texttt{AICMET}.
%
The NLME baseline is fitted study-by-study for each parent compound using one- and two-compartment models with first-order absorption, inter-individual variability, and combined residual error models, selecting the final model by AIC. Metabolites are excluded from this comparison because they lack explicit dosing inputs.
%
All neural models are pretrained on synthetic PK studies and evaluated in context on empirical studies without dataset-specific fine-tuning.
%
NODE-PK extends latent neural ODEs~\citep{rubanova2019latent} with explicit dosing information through an RNN encoder and neural ODE dynamics~\citep{lu2021deep}, but treats individuals independently and does not aggregate study-level information.
%
The Transformer baseline replaces NODE-PK's neural ODE dynamics with a functional-attention decoder, but does not use the full hierarchical mixed-effects representation of \texttt{AICMET}.
%
\texttt{AICMET}~\citep{marin2025amortized} is a hierarchical neural-process baseline that encodes study-level and individual-level effects through latent mixed/random-effect representations and decodes concentration functions with functional attention.
%
Relative to \texttt{AICMET}, PFF learns a direct transport between measures over concentration functions and parameterizes the transport vector field with continuum-attention neural operators rather than a functional-attention decoder.

We additionally evaluate two PFF ablations.
%
\texttt{PFF-white-noise} replaces the GP posterior reference measure used for partially observed individuals with a white-noise reference, isolating the effect of prefix-conditioned functional initialization.
%
\texttt{PFF-std-attn} replaces the continuum-attention neural operator in the vector-field architecture by the standard functional attention mechanism used in \texttt{AICMET} while keeping the flow-matching objective fixed.

\subsection{Metrics}

\paragraph{Forecasting evaluation.}
Forecasting performance is assessed in a leave-one-individual-out setting.
%
For subject $i$, predictions are conditioned on the study context $\mathcal{S}$, dosing inputs $\mathbf{u}^i$, and the first four PK samples $\mathcal{D}^i_{\mathrm{C}}$.
%
The resulting posterior predictive trajectories are evaluated on the remaining samples $\mathcal{D}^i_{\mathrm{T}}$ using log-RMSE.

\paragraph{Generative evaluation.}
We assess sample quality in empirical and synthetic settings.
%
For empirical data, we use leave-one-individual-out evaluation across all real studies, where the model is conditioned on the remaining individuals and generates one trajectory for the held-out individual, yielding a generated set matched in size to the empirical dataset.
%
For the synthetic evaluation, no leave-one-out construction is needed.
%
Instead, each synthetic study provides a study context together with a reference population of target individuals sampled directly from the simulator.
%
Conditioned on the same context, the model generates a matching synthetic population, and the generated trajectories are compared directly against the corresponding oracle trajectories.
After constructing these empirical and synthetic comparison pairs, we quantify discrepancy using two complementary two-sample metrics:
classifier AUC from a two-sample test, where a value of $0.5$ corresponds to perfect agreement, and signature-kernel MMD~\citep{kiraly2019kernels}, where values closer to zero indicate better agreement. Since we use an unbiased finite-sample estimator, small negative MMD estimates can occur.
%
We also report visual predictive checks (VPCs), comparing simulation percentiles against observed concentrations~\citep{holford2005}; implementation details are provided in the Supplementary Material.

\begin{table}[t]
\centering
\caption{\textbf{Generative sample fidelity and ablations.}
AUC values closer to $0.5$ and MMD values closer to zero indicate better agreement between generated and reference trajectories. Values are reported as mean $\pm$ standard error over seeds when available.}
\label{tab:sample_fidelity_ablation}
\smaller
\begin{tabular}{lcc}
\toprule
\textbf{Model} & \textbf{AUC} & \textbf{MMD} \\
\midrule
\multicolumn{3}{l}{\textit{Synthetic simulator-controlled}} \\
\texttt{AICMET} & 0.54 $\pm$ 0.01 & 0.002 $\pm$ 0.001 \\
\texttt{PFF} & \textbf{0.52 $\pm$ 0.01} & \textbf{0.001 $\pm$ 0.001} \\
\texttt{PFF-std-attn} & 0.61 $\pm$ 0.03 & 0.004 $\pm$ 0.004 \\
\texttt{PFF-white-noise} & 0.59 $\pm$ 0.05 & 0.004 $\pm$ 0.001 \\
\addlinespace
\multicolumn{3}{l}{\textit{Empirical leave-one-individual-out}} \\
\texttt{AICMET} & 0.58 $\pm$ 0.01& $-$ \\
\texttt{PFF} & \textbf{0.57 $\pm$ 0.02} & $-$ \\
\bottomrule
\end{tabular}
\end{table}

\subsection{Discussion of Numerical Results}
\label{sec:numerical_results}

Table~\ref{tab:main_results} summarizes forecasting performance.
%
\texttt{PFF} achieves the lowest or competitive log-RMSE across the evaluated compounds in the majority of cases, indicating that the learned posterior predictive operator transfers effectively from simulated pretraining studies to sparse empirical PK datasets.
%
Figure~\ref{fig:prediction_row} illustrates this behavior for partially observed individuals: the model does not merely interpolate observed prefixes, but uses the cohort-level context and individual prefix to infer a distribution over plausible future concentration functions.

Table~\ref{tab:sample_fidelity_ablation} evaluates zero-shot population synthesis and the main PFF ablations.
%
Compared with \texttt{AICMET}, \texttt{PFF} yields AUC values closer to the ideal value of $0.5$ in both empirical and synthetic settings, indicating that its generated trajectories are harder to distinguish from reference trajectories.
%
The synthetic comparison confirms the same trend under simulator-controlled ground truth, while the empirical leave-one-individual-out protocol shows that the improvement transfers to real PK studies.
%
Figure~\ref{fig:vpc_only} provides the corresponding visual predictive checks. The simulated percentiles align closely with the empirical concentration distributions across study cohorts.

The ablations clarify the source of this improvement.
%
Replacing the GP posterior reference with an unconditional white-noise reference worsens MMD, suggesting that prefix-conditioned reference measures provide a useful initialization before transport.
%
Replacing continuum operator attention with standard attention has a larger effect on AUC, making \texttt{PFF-std-attn} nearly indistinguishable from \texttt{AICMET}.
%
Thus, the empirical gain should not be attributed to function-space flow matching alone, but to its combination with GP posterior references and continuum operator attention for irregular observation grids.

Finally, PFF shifts computational cost from repeated study-specific optimization to one up-front pretraining stage.
%
At deployment time, the pretrained model can be applied off-the-shelf to new studies without retraining, whereas classical NLME workflows often require iterative model development, model comparison, and manual diagnostic inspection. 
%
Although PFF generates by integrating a flow ODE ($\sim$100 steps) rather than in a single amortized pass, inference remains cheap in absolute terms. Sampling $100$ virtual individuals for a study takes approximately $3.5$\,s on a single A100 GPU (around $11\times$ the \texttt{AICMET} single-pass baseline), a one-off cost per query that is negligible beside the repeated per-study optimization it replaces.
%
This makes PFF particularly relevant in early-development settings where datasets are small, irregular,  and repeatedly encountered across compounds.
  %
A detailed account of hyperparameters, inference times, calibration diagnostics, and failure cases is provided in the Supplementary Material.

% Finally, PFF shifts computational cost from repeated study-specific optimization to one up-front pretraining stage.
% %
% At deployment time, the pretrained model can be applied off-the-shelf to new studies without retraining, whereas classical NLME workflows often require iterative model development, model comparison, and manual diagnostic inspection.
% %
% This makes PFF particularly relevant in early-development settings where datasets are small, irregular, and repeatedly encountered across compounds.
% %
% A detailed account of hyperparameters, inference times, calibration diagnostics, and failure cases is provided in the Supplementary Material.